% Source: design/adc_topology.scs, sar_comparator.
% @mos sar_comparator M1 NMOS p ip tail vss 3e-6 30e-9
% @mos sar_comparator M2 NMOS n in tail vss 3e-6 30e-9
% @mos sar_comparator M3 NMOS x y p vss 1e-6 30e-9
% @mos sar_comparator M4 NMOS y x n vss 1e-6 30e-9
% @mos sar_comparator M5 PMOS x y vdd vdd 1e-6 30e-9
% @mos sar_comparator M6 PMOS y x vdd vdd 1e-6 30e-9
% @mos sar_comparator M7 NMOS tail clk vss vss 3e-6 30e-9
% @mos sar_comparator MPp PMOS p clk vdd vdd 1e-6 30e-9
% @mos sar_comparator MPn PMOS n clk vdd vdd 1e-6 30e-9
% @mos sar_comparator MPx PMOS x clk vdd vdd 1e-6 30e-9
% @mos sar_comparator MPy PMOS y clk vdd vdd 1e-6 30e-9
% @instance sar_comparator XBX0 sar_inv vdd vss x xb
% @instance sar_comparator XBX1 sar_inv vdd vss xb xd
% @instance sar_comparator XBY0 sar_inv vdd vss y yb
% @instance sar_comparator XBY1 sar_inv vdd vss yb yd
% @instance sar_comparator XLQ sar_nand2 vdd vss xd qb q
% @instance sar_comparator XLQB sar_nand2 vdd vss yd q qb
\begin{circuitikz}
\node[nmos] (m1) at (1,1.5) {};
\node[nmos,xscale=-1] (m2) at (5,1.5) {};
\node[nmos] (m3) at (1,3.25) {};
\node[nmos,xscale=-1] (m4) at (5,3.25) {};
\node[pmos] (m5) at (1,5) {};
\node[pmos,xscale=-1] (m6) at (5,5) {};
\node[nmos] (m7) at (3,-.1) {};
\draw (m5.S) -- (1,5.9) -- (5,5.9) -- (m6.S);
\node[above] at (3,5.9) {$V_{\mathrm{DD}}$};
\draw (m5.D) -- (m3.D) (m6.D) -- (m4.D);
\draw (m3.S) -- (m1.D) (m4.S) -- (m2.D);
\draw (m1.S) -- (1,.65) -- (5,.65) -- (m2.S);
\draw (3,.65) node[adc dot] {} -- (m7.D);
\draw (m7.S) -- (3,-1) node[ground] {};
\draw (m1.G) -- ++(-.45,0) node[left] {$v_p\;(ip)$};
\draw (m2.G) -- ++(.45,0) node[right] {$v_n\;(in)$};
\draw (m7.G) -- ++(-.45,0) node[left] {$\phi_c\;(clk)$};
\foreach \i/\gate in {3/y,5/y}{\draw (m\i.G) -- ++(-.35,0) node[left] {$\gate$};}
\foreach \i/\gate in {4/x,6/x}{\draw (m\i.G) -- ++(.35,0) node[right] {$\gate$};}
\draw (1,4.12) node[adc dot] {} -- ++(-.6,0) node[left] {$x$};
\draw (5,4.12) node[adc dot] {} -- ++(.6,0) node[right] {$y$};
\draw (1,2.35) node[adc dot] {} -- ++(-.6,0) node[left] {$p$};
\draw (5,2.35) node[adc dot] {} -- ++(.6,0) node[right] {$n$};
\node[above] at (3,.65) {$tail$};
\foreach \i in {1,3,5}{\node[right] at (m\i.D) {$M_{\i}$};}
\foreach \i in {2,4,6}{\node[left] at (m\i.D) {$M_{\i}$};}
\node[right] at (m7.D) {$M_7$};
\node at (3,-1.65) {(a) Regenerative core};
% Four reset PMOS devices are shown separately, with connected net labels.
\foreach \xx/\nn in {8.5/p,10.6/n,12.7/x,14.8/y}{
 \node[pmos] (r\nn) at (\xx,5) {};
 \draw (r\nn.S) -- (\xx,5.9);
 \draw (r\nn.D) -- (\xx,4.15) node[below] {$\nn$};
 \draw (r\nn.G) -- ++(-.2,0) node[left] {$\phi_c$};
 \node[right] at (r\nn.D) {$M_{P\nn}$};
}
\draw (8.5,5.9)--(14.8,5.9);
\node[above] at (11.65,5.9) {$V_{\mathrm{DD}}$};
\node at (11.65,3.65) {(b) Precharge: $\phi_c=0$};
% Two CMOS inverters per path; cross-coupled NAND SR hold.
\node[not port] (bx0) at (8.5,2.35) {};
\node[not port] (bx1) at (10.4,2.35) {};
\node[not port] (by0) at (8.5,.35) {};
\node[not port] (by1) at (10.4,.35) {};
\node[nand port] (lq) at (13.6,2.35) {};
\node[nand port] (lqb) at (13.6,.35) {};
\draw (bx0.in)--++(-.35,0) node[left] {$x$};
\draw (by0.in)--++(-.35,0) node[left] {$y$};
\draw (bx0.out)--node[above] {$x_b$}(bx1.in);
\draw (by0.out)--node[above] {$y_b$}(by1.in);
\draw (bx1.out)--(11.7,2.35)|-node[pos=.25,above] {$x_d$}(lq.in 1);
\draw (by1.out)--(11.7,.35)|-node[pos=.25,above] {$y_d$}(lqb.in 1);
\draw (lq.in 2)--++(-.5,0) node[left] {$q_b$};
\draw (lqb.in 2)--++(-.5,0) node[left] {$q$};
\draw[adc signal] (lq.out)--++(.8,0) node[right] {$q$};
\draw[adc signal] (lqb.out)--++(.8,0) node[right] {$q_b$};
\node at (11.5,-.4) {(c) Isolation buffers and SR hold};
\node[align=center] at (11.5,-1.35) {All NMOS bulks: $V_{\mathrm{SS}}$; all PMOS bulks: $V_{\mathrm{DD}}$.\\Equal net labels denote electrical connections.};
\end{circuitikz}
